% GENERATED FILE. Edit canonical lessons, then run npm run generate:books.
% canonical-source-hash: fnv1a64:95b3cfa6ad9917d6
% canonical-lessons: GU-C470-cahero, GU-C470-papan, GU-C470-cotlo, GU-C470-lat, GU-C470-pathi, GU-R470-review-463-470

\chapter{Face and Hair}
\label{ch:gu-face-and-hair}
\hlchaptermodality{470}

\begin{chapteropening}
\textbf{By the end of this chapter:} \emph{I can describe a face and a hairstyle.}

Complete the last lesson of chapter 470: I can describe a face and a hairstyle.
\end{chapteropening}

\section[\texorpdfstring{\gu{ચહેરો} (cahero)}{ચહેરો (cahero)}]{\gu{ચહેરો} (cahero) — a face}

\label{lesson:GU-C470-cahero}



\begin{quote}
Before the new one: say the Gujarati for sluggish, then the Gujarati for steady effort.
\end{quote}

\subsection*{You'll want to know: \gu{ચહેરો}}
\textbf{\gu{ચહેરો}} — \emph{cahero} — \textquotedblleft{}a face\textquotedblright{}.

\begin{quote}
\gu{મીરાનો} \gu{ચહેરો} \gu{ખુશીથી} \gu{ચમકતો} \gu{હતો}.
\end{quote}

\begin{quote}
\emph{Mīrāno cahero khuśīthī camakto hato.}
\end{quote}

\begin{quote}
Mira's face was shining with happiness.
\end{quote}

A masculine noun.

\begin{cousinweb}
\textbf{\gu{ચહેરો}} (\emph{cahero}) is usually traced to Persian \emph{čihra}, \textquotedblleft{}face, appearance\textquotedblright{}.
\end{cousinweb}

\begin{grammarlens}[title={What you've built}]
One more word for describing what you lived and felt.
\end{grammarlens}

\subsection*{Your turn}
\begin{itemize}
\raggedright
  \item \emph{Say it:} \emph{cahero}
  \item \emph{Say it:} \emph{cahero}, once more
  \item \emph{Say it:} \emph{Mīrāno cahero khuśīthī camakto hato.} — Mira's face was shining with happiness.
  \item \emph{Recall:} say the Gujarati for sluggish, then the Gujarati for steady effort, then say \emph{cahero} again
\end{itemize}

\begin{tcolorbox}[breakable,colback=teal!4,colframe=teal!35!black,title={Before you move on}]
What does \gu{ચહેરો} mean? (\textquotedblleft{}A face\textquotedblright{}.) Say it once more.
\end{tcolorbox}
\section[\texorpdfstring{\gu{પાંપણ} (pā̃paṇ)}{પાંપણ (pā̃paṇ)}]{\gu{પાંપણ} (pā̃paṇ) — an eyelash}

\label{lesson:GU-C470-papan}



\begin{quote}
Before the new one: say the Gujarati for steady effort, then the Gujarati for a face.
\end{quote}

\subsection*{You'll want to know: \gu{પાંપણ}}
\textbf{\gu{પાંપણ}} — \emph{pā̃paṇ} — \textquotedblleft{}an eyelash\textquotedblright{}.

\begin{quote}
\gu{બાળકની} \gu{પાંપણ} \gu{લાંબી} \gu{છે}.
\end{quote}

\begin{quote}
\emph{bāḷaknī pā̃paṇ lā̃bī che.}
\end{quote}

\begin{quote}
The baby has long eyelashes.
\end{quote}

A feminine noun.

\begin{cousinweb}
\textbf{\gu{પાંપણ}} (\emph{pā̃paṇ}) is inherited from Sanskrit \emph{pakṣman}, \textquotedblleft{}eyelash\textquotedblright{}; CDIAL lists the Gujarati word (CDIAL 7638).
\end{cousinweb}

\begin{grammarlens}[title={What you've built}]
One more word for describing what you lived and felt.
\end{grammarlens}

\subsection*{Your turn}
\begin{itemize}
\raggedright
  \item \emph{Say it:} \emph{pā̃paṇ}
  \item \emph{Say it:} \emph{pā̃paṇ}, once more
  \item \emph{Say it:} \emph{bāḷaknī pā̃paṇ lā̃bī che.} — The baby has long eyelashes.
  \item \emph{Recall:} say the Gujarati for steady effort, then the Gujarati for a face, then say \emph{pā̃paṇ} again
\end{itemize}

\begin{tcolorbox}[breakable,colback=teal!4,colframe=teal!35!black,title={Before you move on}]
What does \gu{પાંપણ} mean? (\textquotedblleft{}An eyelash\textquotedblright{}.) Say it once more.
\end{tcolorbox}
\section[\texorpdfstring{\gu{ચોટલો} (coṭlo)}{ચોટલો (coṭlo)}]{\gu{ચોટલો} (coṭlo) — a plait of hair}

\label{lesson:GU-C470-cotlo}



\begin{quote}
Before the new one: say the Gujarati for a face, then the Gujarati for an eyelash.
\end{quote}

\subsection*{You'll want to know: \gu{ચોટલો}}
\textbf{\gu{ચોટલો}} — \emph{coṭlo} — \textquotedblleft{}a plait of hair\textquotedblright{}.

\begin{quote}
\gu{દાદી} \gu{રોજ} \gu{સવારે} \gu{મીરાનો} \gu{ચોટલો} \gu{ગૂંથે} \gu{છે}.
\end{quote}

\begin{quote}
\emph{dādī roj savāre Mīrāno coṭlo gū̃the che.}
\end{quote}

\begin{quote}
Every morning my grandmother plaits Mira's hair.
\end{quote}

A masculine noun.

\begin{cousinweb}
\textbf{\gu{ચોટલો}} (\emph{coṭlo}) is inherited: CDIAL lists the Gujarati word, \textquotedblleft{}hair of head\textquotedblright{}, under Sanskrit \emph{cūḍā}, \textquotedblleft{}topknot on the head\textquotedblright{} (CDIAL 4883).
\end{cousinweb}

\begin{grammarlens}[title={What you've built}]
One more word for describing what you lived and felt.
\end{grammarlens}

\subsection*{Your turn}
\begin{itemize}
\raggedright
  \item \emph{Say it:} \emph{coṭlo}
  \item \emph{Say it:} \emph{coṭlo}, once more
  \item \emph{Say it:} \emph{dādī roj savāre Mīrāno coṭlo gū̃the che.} — Every morning my grandmother plaits Mira's hair.
  \item \emph{Recall:} say the Gujarati for a face, then the Gujarati for an eyelash, then say \emph{coṭlo} again
\end{itemize}

\begin{tcolorbox}[breakable,colback=teal!4,colframe=teal!35!black,title={Before you move on}]
What does \gu{ચોટલો} mean? (\textquotedblleft{}A plait of hair\textquotedblright{}.) Say it once more.
\end{tcolorbox}
\section[\texorpdfstring{\gu{લટ} (laṭ)}{લટ (laṭ)}]{\gu{લટ} (laṭ) — a lock of hair}

\label{lesson:GU-C470-lat}



\begin{quote}
Before the new one: say the Gujarati for an eyelash, then the Gujarati for a plait.
\end{quote}

\subsection*{You'll want to know: \gu{લટ}}
\textbf{\gu{લટ}} — \emph{laṭ} — \textquotedblleft{}a lock of hair\textquotedblright{}.

\begin{quote}
\gu{પવનથી} \gu{તેના} \gu{કપાળ} \gu{પર} \gu{એક} \gu{લટ} \gu{આવી} \gu{ગઈ}.
\end{quote}

\begin{quote}
\emph{pavanthī tenā kapāḷ par ek laṭ āvī gaī.}
\end{quote}

\begin{quote}
The wind blew a lock of hair onto her forehead.
\end{quote}

A feminine noun.

\begin{cousinweb}
\textbf{\gu{લટ}} (\emph{laṭ}) probably goes back to a reconstructed \emph{*laṭṭa-}, \textquotedblleft{}hanging, tangled\textquotedblright{}, under which CDIAL lists the Gujarati \emph{laṭiyā̃}, \textquotedblleft{}matted locks\textquotedblright{} (CDIAL 10918).
\end{cousinweb}

\begin{grammarlens}[title={What you've built}]
One more word for describing what you lived and felt.
\end{grammarlens}

\subsection*{Your turn}
\begin{itemize}
\raggedright
  \item \emph{Say it:} \emph{laṭ}
  \item \emph{Say it:} \emph{laṭ}, once more
  \item \emph{Say it:} \emph{pavanthī tenā kapāḷ par ek laṭ āvī gaī.} — The wind blew a lock of hair onto her forehead.
  \item \emph{Recall:} say the Gujarati for an eyelash, then the Gujarati for a plait, then say \emph{laṭ} again
\end{itemize}

\begin{tcolorbox}[breakable,colback=teal!4,colframe=teal!35!black,title={Before you move on}]
What does \gu{લટ} mean? (\textquotedblleft{}A lock of hair\textquotedblright{}.) Say it once more.
\end{tcolorbox}
\section[\texorpdfstring{\gu{પાંથી} (pā̃thī)}{પાંથી (pā̃thī)}]{\gu{પાંથી} (pā̃thī) — the parting of the hair}

\label{lesson:GU-C470-pathi}



\begin{quote}
Before the new one: say the Gujarati for a plait, then the Gujarati for a lock of hair.
\end{quote}

\subsection*{You'll want to know: \gu{પાંથી}}
\textbf{\gu{પાંથી}} — \emph{pā̃thī} — \textquotedblleft{}the parting of the hair\textquotedblright{}.

\begin{quote}
\gu{મીરાની} \gu{પાંથી} \gu{વચ્ચે} \gu{છે}.
\end{quote}

\begin{quote}
\emph{Mīrānī pā̃thī vacce che.}
\end{quote}

\begin{quote}
Mira parts her hair in the middle.
\end{quote}

A feminine noun.

\begin{cousinweb}
\textbf{\gu{પાંથી}} (\emph{pā̃thī}) is inherited from Sanskrit \emph{pánthā}, \textquotedblleft{}path\textquotedblright{}; CDIAL lists the Gujarati word, \textquotedblleft{}parting of the hair\textquotedblright{} (CDIAL 7785): a little path through the hair.
\end{cousinweb}

\begin{grammarlens}[title={What you've built}]
One more word for describing what you lived and felt.
\end{grammarlens}

\subsection*{Your turn}
\begin{itemize}
\raggedright
  \item \emph{Say it:} \emph{pā̃thī}
  \item \emph{Say it:} \emph{pā̃thī}, once more
  \item \emph{Say it:} \emph{Mīrānī pā̃thī vacce che.} — Mira parts her hair in the middle.
  \item \emph{Recall:} say the Gujarati for a plait, then the Gujarati for a lock of hair, then say \emph{pā̃thī} again
\end{itemize}

\begin{tcolorbox}[breakable,colback=teal!4,colframe=teal!35!black,title={Before you move on}]
What does \gu{પાંથી} mean? (\textquotedblleft{}The parting of the hair\textquotedblright{}.) Say it once more.
\end{tcolorbox}
\section[(dialogue)]{Looking back: from Chatty and Quiet to Face and Hair}

\label{lesson:GU-R470-review-463-470}



\begin{quote}
Say the Gujarati for straightforward and of few words.
\end{quote}

\subsection*{Your turn}
Say these aloud:

\begin{itemize}
\raggedright
  \item straightforward, of few words, talkative, to babble, cheerful
  \item discreet, good sense, well-behaved, composed, to heed
  \item strict, fussy, insolent, egotism, to get used to
  \item cowardly, fearless, valour, grit, to wobble
  \item behaviour, respectful manners, airs, moral character, to be ashamed
  \item untruthful, a tall tale, hypocrisy, boasting, to beguile
  \item odd, whimsical, obsession, sluggish, steady effort
  \item a face, an eyelash, a plait, a lock of hair, the parting of the hair
\end{itemize}

\begin{tcolorbox}[breakable,colback=teal!4,colframe=teal!35!black,title={Before you move on}]
Straightforward? (\textbf{\gu{સરળ}}, \emph{saraḷ}.) Of few words? (\textbf{\gu{ઓછાબોલું}}, \emph{ochābolũ}.) Talkative? (\textbf{\gu{વાચાળ}}, \emph{vācāḷ}.) To babble? (\textbf{\gu{બકવું}}, \emph{bakvũ}.) Cheerful? (\textbf{\gu{હસમુખું}}, \emph{hasmukhũ}.) Discreet? (\textbf{\gu{શાણું}}, \emph{śāṇũ}.) Good sense? (\textbf{\gu{વિવેક}}, \emph{vivek}.) Well-behaved? (\textbf{\gu{ડાહ્યું}}, \emph{ḍāhyũ}.) Composed? (\textbf{\gu{ઠાવકું}}, \emph{ṭhāvkũ}.) To heed? (\textbf{\gu{ગાંઠવું}}, \emph{gā̃ṭhvũ}.) Strict? (\textbf{\gu{સખત}}, \emph{sakhat}.) Fussy? (\textbf{\gu{ચીકણું}}, \emph{cīkṇũ}.) Insolent? (\textbf{\gu{ઉદ્ધત}}, \emph{uddhat}.) Egotism? (\textbf{\gu{અહંકાર}}, \emph{ahankār}.) To get used to? (\textbf{\gu{ટેવાવું}}, \emph{ṭevāvũ}.) Cowardly? (\textbf{\gu{કાયર}}, \emph{kāyar}.) Fearless? (\textbf{\gu{નીડર}}, \emph{nīḍar}.) Valour? (\textbf{\gu{શૌર્ય}}, \emph{śaurya}.) Grit? (\textbf{\gu{ખમીર}}, \emph{khamīr}.) To wobble? (\textbf{\gu{ડગમગવું}}, \emph{ḍagmagvũ}.) Behaviour? (\textbf{\gu{વર્તન}}, \emph{vartan}.) Respectful manners? (\textbf{\gu{અદબ}}, \emph{adab}.) Airs? (\textbf{\gu{નખરાં}}, \emph{nakhrā̃}.) Moral character? (\textbf{\gu{ચારિત્ર્ય}}, \emph{cāritrya}.) To be ashamed? (\textbf{\gu{લાજવું}}, \emph{lājvũ}.) Untruthful? (\textbf{\gu{જૂઠું}}, \emph{jūṭhũ}.) A tall tale? (\textbf{\gu{ગપ્પું}}, \emph{gappũ}.) Hypocrisy? (\textbf{\gu{દંભ}}, \emph{dambh}.) Boasting? (\textbf{\gu{શેખી}}, \emph{śekhī}.) To beguile? (\textbf{\gu{ભોળવવું}}, \emph{bhoḷavvũ}.) Odd? (\textbf{\gu{વિચિત્ર}}, \emph{vicitra}.) Whimsical? (\textbf{\gu{તરંગી}}, \emph{tarangī}.) Obsession? (\textbf{\gu{ઝનૂન}}, \emph{jhanūn}.) Sluggish? (\textbf{\gu{સુસ્ત}}, \emph{sust}.) Steady effort? (\textbf{\gu{ઉદ્યમ}}, \emph{udyam}.) A face? (\textbf{\gu{ચહેરો}}, \emph{cahero}.) An eyelash? (\textbf{\gu{પાંપણ}}, \emph{pā̃paṇ}.) A plait? (\textbf{\gu{ચોટલો}}, \emph{coṭlo}.) A lock of hair? (\textbf{\gu{લટ}}, \emph{laṭ}.) The parting of the hair? (\textbf{\gu{પાંથી}}, \emph{pā̃thī}.)
\end{tcolorbox}
